\section{Votación Electrónica y Propiedades de Seguridad}
% =============================================================================

Un sistema de votación electrónica se define formalmente como el conjunto de procedimientos criptográficos, protocolos de comunicación, arquitecturas de software e infraestructuras de cómputo que permiten la acreditación de electores, la emisión confidencial del sufragio, su registro inmutable, su transmisión segura y el cómputo final de resultados mediante medios digitales (Gritzalis, 2002; Volkamer, 2009). 

A diferencia de los sistemas de información transaccionales convencionales ---tales como las plataformas bancarias, el comercio electrónico o la gestión gubernamental en línea---, el diseño de un sistema electoral digital enfrenta requerimientos de seguridad estrictamente antagónicos: debe garantizar simultáneamente la autenticidad e identificación inequívoca del ciudadano en la fase de acceso, y la desvinculación estricta y permanente entre dicha identidad y el contenido de su papeleta en las fases de custodia y escrutinio (Chaum, 1981; Chaum, 1982; Benaloh, 1987).

\subsection{Evolución Histórica y Enfoques Criptográficos en Sistemas Electorales}

Desde la formulación de los primeros esquemas criptográficos en la década de 1980, la literatura especializada ha explorado tres grandes paradigmas para resolver la tensión entre anonimato y auditabilidad:

\begin{enumerate}
	\item \textbf{Redes de Mezcla (\textit{Mix-Nets}):} Propuestas originalmente por David Chaum (1981), se basan en una secuencia de servidores independientes denominados \textit{mixes}. Cada servidor recibe un lote de boletas cifradas con capas sucesivas de criptografía de clave pública, aplica una permutación aleatoria secreta a los elementos y remueve una capa de cifrado antes de retransmitirlos al siguiente nodo. Si al menos un servidor de la cadena es honesto, se destruye la correlación temporal y espacial entre el emisor y el mensaje saliente. Variantes modernas incorporan pruebas de conocimiento cero de permutación correcta (\textit{verifiable shuffles}) (Neff, 2001; Cortier et al., 2014).
	\item \textbf{Cifrado Homomórfico y Pruebas de Conocimiento Cero ($ZKP$):} Introducido en el ámbito electoral por Benaloh (1987) y formalizado por Cramer, Gennaro y Schoenmakers (1997), este enfoque permite computar operaciones algebraicas sobre los votos directamente en su forma cifrada. Empleando criptosistemas aditivamente homomórficos (como Paillier o ElGamal exponencial), las autoridades multiplican los textos cifrados para obtener el cifrado de la suma de los votos, descifrando únicamente el resultado final consolidado sin abrir jamás una boleta individual.
	\item \textbf{Firmas Ciegas (\textit{Blind Signatures}):} Diseñadas por Chaum (1988) y aplicadas al sufragio por Fujioka, Okamoto y Ohta (1992) en el seminal protocolo FOO. El votante introduce su papeleta dentro de un ``sobre ciego digital'' mediante un factor de enmascaramiento aleatorio; la autoridad firma el sobre opaco certificando que el elector está empadronado, y posteriormente el votante retira la máscara, obteniendo una credencial válidamente firmada que puede depositar anónimamente en el ánfora a través de un canal desvinculado.
\end{enumerate}

\subsection{Taxonomía Formal de las Propiedades de Seguridad Electoral}

La doctrina de seguridad de la información electoral clasifica los requisitos que debe satisfacer cualquier protocolo comicial en cuatro dimensiones fundamentales (Volkamer, 2009; Bernhard et al., 2012; Cortier et al., 2014):

\subsubsection{Confidencialidad y Privacidad del Sufragio (\textit{Ballot Secrecy})}
El secreto del voto constituye un principio democrático consagrado internacionalmente. Desde la teoría criptográfica, la confidencialidad exige que ningún actor ---incluyendo los administradores de la red, los proveedores de infraestructura y las propias autoridades electorales coludidas--- sea capaz de asociar la identidad de un votante con la opción política seleccionada (Bernhard et al., 2012).

En la literatura avanzada se distinguen tres niveles crecientes de privacidad:
\begin{itemize}
	\item \textbf{Secreto Simple (\textit{Ballot Privacy}):} Ningún observador pasivo o activo del tráfico de red puede deducir la preferencia contenida en la boleta con una ventaja estadística superior a la de una estimación aleatoria.
	\item \textbf{Libertad de Recibo (\textit{Receipt-Freeness}):} Formalizada por Benaloh y Tuinstra (1994) y Delaune, Kremer y Ryan (2009), establece que el sistema no debe entregar al votante ningún elemento o comprobante que le permita demostrar fehacientemente a un tercero por quién votó. Esta propiedad neutraliza de raíz los esquemas comerciales de compra y venta de votos (\textit{vote buying}).
	\item \textbf{Resistencia a la Coacción (\textit{Coercion Resistance}):} Formulada por Juels, Catalano y Jakobsson (2010), constituye la garantía más exigente: asegura que el votante pueda resistir la extorsión o intimidación directa de un coactor, incluso en escenarios donde el coactor le ordene compartir sus credenciales secretas, le instruya qué votar o presencie físicamente la pantalla del dispositivo durante la emisión.
\end{itemize}

\subsubsection{Verificabilidad de Extremo a Extremo (\textit{End-to-End Verifiability, E2E-V})}
A diferencia de los sistemas de cómputo cerrados en los que la sociedad debe depositar un ``acto de fe ciega'' en los ingenieros de software o en los fabricantes de hardware, los protocolos criptográficos modernos implementan el principio de \textbf{Verificabilidad de Extremo a Extremo} (Chaum, 2004; Adida, 2008; Ryan et al., 2009). Formalizada por Cortier, Delaune e Issarny (2014), la verificabilidad E2E-V exige el cumplimiento concurrente de tres propiedades matemáticas auditables:

\begin{enumerate}
	\item \textbf{Emitido según la intención (\textit{Cast-as-Intended}):} El votante tiene la certeza técnica de que el criptograma generado en su terminal refleja fielmente la opción que marcó en pantalla, sin que un programa malicioso (\textit{malware}) en el navegador o en la cabina haya sustituido su decisión. Esto se resuelve comúnmente mediante el mecanismo interactivo de auditoría \textit{Benaloh Challenge} (\textit{cast-or-spoil}) (Benaloh, 1987) o mediante esquemas de comprobación visual criptográfica (Chaum, 2004).
	\item \textbf{Registrado según lo emitido (\textit{Recorded-as-Cast}):} El votante recibe un comprobante público o código de seguimiento determinista (como un hash inmutable) que le permite verificar en un tablero o boletín de acceso abierto (\textit{public bulletin board}) que su papeleta cifrada fue recibida e incorporada formalmente en la urna digital sin alteraciones ni omisiones.
	\item \textbf{Contabilizado según lo registrado (\textit{Tallied-as-Recorded}):} Cualquier ciudadano, fiscal partidario u observador internacional puede auditar matemáticamente el escrutinio final a partir del conjunto íntegro de boletas publicadas en el boletín, constatando que el cómputo oficial deriva de forma exacta y no corrompida de las papeletas archivadas en la urna, sin necesidad de confiar en la honestidad de las autoridades.
\end{enumerate}

\subsubsection{Elegibilidad, Unicidad del Sufragio y No Repudio}
\begin{itemize}
	\item \textbf{Elegibilidad (\textit{Eligibility}):} Únicamente los ciudadanos legalmente empadronados y habilitados por la Autoridad de Registro pueden inyectar un voto válido en el sistema (Sampigethaya y Poovendran, 2006).
	\item \textbf{Unicidad (\textit{One-Person One-Vote / Unreusability}):} Cada elector habilitado puede emitir exactamente un voto admisible. Cualquier intento deliberado o accidental de depositar papeletas adicionales debe ser detectado y neutralizado en tiempo real por el servidor de recepción (Fujioka et al., 1992).
	\item \textbf{No Repudio e Integridad de la Emisión:} El sufragio legítimo debe incorporar una vinculación criptográfica irrefutable (mediante firma digital) que impida al elector desconocer la autoría de una emisión válida, protegiendo al sistema ante impugnaciones fraudulentas posteriores (Cramer et al., 1997).
\end{itemize}

\subsubsection{Robustez y Tolerancia a Fallos (\textit{Robustness})}
El protocolo debe ser capaz de culminar el proceso comicial y generar un resultado incontrovertible incluso ante la presencia de votantes maliciosos que envíen paquetes corrompidos, la caída transitoria de enlaces de telecomunicaciones o la deserción de un subconjunto de autoridades electorales (Cramer et al., 1997; Chevallier-Mames et al., 2010).

\subsection{La Paradoja de la Votación Electrónica frente a la Banca Digital}

Diversos autores (Saltman, 2006; Volkamer, 2009; Bernhard et al., 2017) han destacado la singular complejidad conceptual del voto digital frente a otras aplicaciones críticas de la computación distribuida:

\begin{quote}
	\small\textit{``En una transacción bancaria o de comercio electrónico, la seguridad se basa en la correlación perfecta y auditable: el cliente exige un recibo firmado que documente con la máxima precisión su identidad, la fecha exacta y el monto transferido, de modo que dicho comprobante pueda emplearse como prueba ante una disputa legal. En la votación democrática, por el contrario, la ley electoral prohíbe de manera terminante entregarle al ciudadano un recibo que exponga su elección en claro, pues dicho comprobante se transformaría en el instrumento predilecto para la compra de votos, el clientelismo y la coacción patronal o política.''} (Adaptado de Saltman, 2006; Volkamer, 2009).
\end{quote}

Esta aparente contradicción ---demostrar matemáticamente que un voto fue fielmente registrado y escrutado, sin entregar ningún comprobante que delate el contenido de la elección--- es el desafío central que la criptografía asimétrica moderna resuelve mediante el uso de funciones unidireccionales con trampa, compromisos criptográficos y resúmenes hash públicos deterministas.

\subsection{Lecciones de Vulnerabilidades en Protocolos Abiertos (Helios y Belenios)}

La viabilidad práctica de los protocolos verificables de extremo a extremo quedó demostrada con el desarrollo de plataformas de código abierto como \textbf{Helios Voting} (Adida, 2008), utilizada extensamente en elecciones académicas internacionales (incluyendo la Universidad de Lovaina y la Universidad de Princeton), y su sucesora europea \textbf{Belenios} (Cortier et al., 2019).

No obstante, la literatura contemporánea ha evidenciado que la robustez matemática aislada de las primitivas criptográficas no garantiza por sí misma la seguridad del sistema si existen deficiencias en el diseño lógico de las interacciones. Baloglu, Bursuc, Mauw y Pang (2021) aplicaron técnicas avanzadas de verificación formal automatizada sobre Helios y Belenios, identificando vulnerabilidades sutiles que habían eludido más de una década de auditorías manuales:
\begin{itemize}
	\item Ataques de sustitución y maleabilidad de textos cifrados (\textit{ballot-stuffing / replay}), en los que un adversario interceptaba boletas válidas y las reinyectaba bajo identificadores artificiales alterando la distribución estadística del conteo.
	\item Fugas de privacidad derivadas de la sincronización de credenciales institucionales compartidas entre la autoridad de autenticación y los servidores de depósito.
\end{itemize}

Los hallazgos de Baloglu et al. (2021) establecieron una conclusión metodológica determinante para la ciencia de la computación: los protocolos de votación no deben ser validados únicamente mediante inspección empírica o pruebas convencionales de software, sino a través de \textbf{modelos de verificación formal basados en lógica matemática y reescritura de términos}, sometidos a las condiciones adversariales más hostiles posibles.

\subsection{El Principio de la Función Unidireccional con Trampa}

Desde la perspectiva criptográfica, la seguridad de cualquier esquema de votación electrónica descansa sobre el concepto de \textit{funciones unidireccionales con trampa} (\textit{trapdoor one-way functions}) (Diffie y Hellman, 1976). Una función $f: X \to Y$ se considera unidireccional si calcular la imagen $y = f(x)$ es computacionalmente eficiente (tiempo polinomial), pero invertir la operación para hallar una preimagen $x \in f^{-1}(y)$ es intratable para cualquier algoritmo clásico eficiente, a menos que se disponga de una información privilegiada adicional denominada \textit{trampa secreta} (\textit{trapdoor}).

\begin{figure}[htbp]
	\centering
	\begin{tikzpicture}[
		scale=0.9,
		node distance=1.1cm,
		cbox/.style={rectangle, draw=black, fill=white, line width=0.8pt, rounded corners=1.5pt, align=center, font=\footnotesize, inner sep=5pt}
	]
		\node (in) [cbox, text width=2.8cm] {\textbf{Mensaje / Voto}\\$x \in X$};
		\node (proc) [cbox, right=of in, text width=3.4cm] {\textbf{Cálculo Criptográfico}\\$y = f(x)$ \ (Polinomial)};
		\node (out) [cbox, right=of proc, text width=2.8cm] {\textbf{Criptograma / Boleta}\\$y \in Y$};
		
		\draw[-{Latex[length=2mm, width=1.5mm]}, line width=0.8pt] (in) -- (proc);
		\draw[-{Latex[length=2mm, width=1.5mm]}, line width=0.8pt] (proc) -- (out);
		
		\draw[-{Latex[length=2mm, width=1.5mm]}, line width=0.7pt, dashed] (out.north) to[bend right=40] 
		node[midway, above, font=\scriptsize\itshape] {Inversión intratable sin trampa} (in.north);
		
		\node (trap) [cbox, below=0.85cm of proc, text width=4.0cm] {\textbf{Trampa Secreta (Clave Privada)}\\Permite calcular $x = f^{-1}(y)$ en tiempo eficiente};
		\draw[-{Latex[length=2mm, width=1.5mm]}, line width=0.8pt] (trap) -- (proc);
	\end{tikzpicture}
	\caption{Principio estructural de una función unidireccional con trampa.}
	\label{fig:trapdoor_func}
	\par\vspace{3pt}{\footnotesize\textit{Fuente: Adaptado de Diffie y Hellman (1976).}}
\end{figure}

En un protocolo de votación, la papeleta es transformada mediante la clave pública electoral en un criptograma indescifrable para el canal, mientras que la autoridad de cómputo utiliza la clave privada como trampa matemática para la apertura legítima al cierre de la jornada electoral.

% =============================================================================